\paragraph{One box can require two source occurrences.}
Linear covers do not identify occurrence count with source-box rank,
even in a source-separable interval above a globally minimal core baseline.

\begin{proposition}[A four-variable failure of maximal-type attainment]
\label{core:prop:one-box-two-copies}
For every $K\geq3$ there is a source-separable core operation
$F(X,x,y,z)$ whose distinguished source family is one normalized owner
box, but whose least number of $X$-occurrences is two, both in the
primitive signature and with union named. Its neutral baseline is
globally minimal among core operations. At $X$, it has no greatest
actual full type.
\end{proposition}
\begin{proof}
Fix distinct owners $a,b,c$, write $m$ for the two-out-of-three
coalition predicate on $V=\{x,y,z\}$, and let $C$ be the greatest
profile completion with
\[
 f_a=f_b=m,\qquad
 f_t=x\land y\land z\quad(t\notin\{a,b\}).
\]
Two two-element subsets of $V$ intersect, so these profiles are
admissible. Choose passive core helpers with profiles
\[
 \begin{array}{c|cccc}
     &a&b&c&t\notin\{a,b,c\}\\\hline
 H_a&m&x\land y&x\land y&x\land y\land z\\
 H_b&x\land y&m&x\land y&x\land y\land z
 \end{array}
\]
by profile synthesis. Their admissibility says that any two-element
coalition meets $\{x,y\}$. Put
\[
 S_a=(X\star_c H_a)\star_b C,\qquad
 S_b=(X\star_c H_b)\star_a C,\qquad
 F=S_a\star_c S_b.
\]
Each spine lies between $C$ and $C\cup X$. For $S_a$, the $a$-singleton
stratum of $C$ makes $H_a=\{a\}$; its $b$-singleton stratum is
restored by the final gate. Every other singleton stratum forces all
passive inputs and the helper to that singleton. Otherwise $C$ is the
whole passive union, containing $H_a$; the final gate cannot lose $C$
by absorbing at $b$, since $H_a=\{b\}$ requires $x=y=\{b\}$,
which would make $C=\{b\}$. The argument for $S_b$ is symmetric.
Thus either spine is $\{c\}$ only when $C=\{c\}$, when both are.
Consequently their final $c$-choice is their union, and $F\in\mathcal I(C)$.
At a distinguished-source probe the spines give
\[
 P=\neg(x_cy_c)\land\bigl(\neg m(E_b)\lor\neg m(E_a)\bigr)
   =\neg(x_cy_c),
\]
since disjoint atomic sets $E_a,E_b$ cannot both have two elements.
Thus the canonical full factors at $X$ are
\[
 B_c(E)=[\{x,y\}\nsubseteq E],\qquad B_t=1\quad(t\ne c).
\]
The passive predicates are those of $C$, which factor ownerwise as
$\bigwedge_t\neg f_t(E_t)$ on proper source probes. Hence $F$ is
source-separable, and $b(P)=1$.

The baseline is globally minimal. In fact it is forced by its positive
profile bits. On a singleton stratum this is immediate. Otherwise
$|E_a|,|E_b|\leq1$. Given $k\in C(w)=\bigcup_iw_i$, choose
$j$ with $k\in w_j$. Choose two-element coalitions $A_a,A_b$ with
$A_a\cap A_b=\{j\}$, avoiding $E_b$ in $A_a$ when $b\ne k$ and
$E_a$ in $A_b$ when $a\ne k$. Such a choice exists: the two forbidden
atomic sets are disjoint and have at most one element, neither equal
to $j$ when its avoidance is required. Take $V$ at all other owners.
These coalitions certify $k\in\mathcal R_f(w)$, so the universal
lower-bound argument in \cref{core:thm:read-once-minimum} forces $k$
in every core operation retaining these profile bits. In particular
$H\subseteq C$ implies $C\subseteq H$.

The canonical full type $B$ fails condition~(2) of
\cref{core:thm:profile-interval-types}: its active coalition
$\{x,y\}$ at $c$ and the inactive winning coalitions
$\{x,z\}$ at $a$ and $\{y,z\}$ at $b$ have empty common intersection.
It is therefore nonactual, and
\cref{core:cor:profile-interval-greatest-type} excludes a greatest
actual type. A one-box cover would have to use precisely this canonical
type, as pure probes recover each factor. The exact source-count
criterion thus excludes one $X$ with union named. The displayed primitive
two-copy realizer proves optimality in both signatures.
\end{proof}

\paragraph{Exact types can force unbounded owner alternation.}
Although a linear cover removes repeated distinguished leaves, its
spines need not have bounded owner alternation, even at fixed $K$.

\begin{proposition}[Ordered reset witnesses]
\label{core:prop:ordered-reset-witnesses}
Fix distinct owners $a,b,c$. For every $m\geq1$ there is a
source-separable core operation on $X,v_1,\ldots,v_m$, with one-box
distinguished source family and a one-$X$ realizer with union named,
such that every primitive one-$X$ spine minorant attaining its canonical
full type has at least $m$ gates.
\end{proposition}
\begin{proof}
Let $r_i$ alternate between $a$ and $b$, and take the left-associated
spine
\[
 T_m=X\star_{r_1}v_1\star_{r_2}\cdots\star_{r_m}v_m,
 \qquad F_m=T_m+\partial_XT_m.
\]
Its canonical full $X$-type is
\[
 B_t(E)=\bigwedge_{i:r_i=t}[v_i\notin E].
\]
The original word is primitive read-once, so every source coefficient
is a single owner box. Neutral restriction at $X$ preserves that
factorization for the passive sources, and sink completion preserves
the $X$ predicate. Thus $F_m$ is separable; $T_m$ is a spine minorant
attaining $B$ in $m$ gates.

Consider any primitive one-$X$ spine $S\subseteq F_m$ with this type.
Every gate on it is owned by $a$ or $b$: a nonempty foreign-owned
helper has positive top singleton profile, contradicting its
identically true factor $B_t$. For each $i$, let $L_i$ be the last
$r_i$-owned gate whose helper's $r_i$-profile is positive at
$\{v_i\}$. Such a gate exists because $B_{r_i}(\{v_i\})=0$.

For $i<j$ with $r_i\ne r_j$, set $v_i=\{r_i\}$,
$v_j=\{r_j\}$, and every other passive input and $X$ to $\{c\}$.
In $T_m$ the reset at $j$ erases $r_i$, and no later input introduces
it. Its neutral residual also omits $r_i$: when querying that winner,
the reset at $i$ already makes the original and neutralized prefixes
equal. Thus $F_m$ omits $r_i$ on this row.
In $S$, the only possible absorbing helpers are detected by their
$r_i$-profile at $\{v_i\}$ or their $r_j$-profile at $\{v_j\}$.
If $L_i>L_j$, its last absorbing helper resets the spine to $r_i$;
every subsequent gate preserves that winner, a contradiction.
Therefore $L_i<L_j$. In particular adjacent alternation forces
$L_1<\cdots<L_m$, giving the claimed lower bound.
\end{proof}

\paragraph{One conditional capture can force the whole order.}
With two spare owners, the ordered witnesses also apply to a single
capture query, with every other marked factor unrestricted.

\begin{proposition}[Exact depth for a single capture]
\label{core:prop:single-capture-depth}
Fix distinct owners $a,b,c,d$ and an odd $m\geq1$. Let $r_i=a$
for odd $i$ and $r_i=b$ for even $i$, and put
\[
 T_m=X\star_{r_1}v_1\star_{r_2}\cdots\star_{r_m}v_m,
 \qquad C_m=\partial_XT_m.
\]
Every one-$X$ spine $S\subseteq C_m\cup X$ whose marked capture
profile satisfies $\psi_a(\{v_1\})=1$ has at least $m$ choice gates
on its marked path. At least $(m+1)/2$ are $a$-owned and have
pairwise distinct helper profiles. Union attachments are allowed.
The depth bound is attained by $T_m$, also with positivity at the probe
\[
 p_a=\{v_i:i\text{ even}\},\qquad
 p_b=\{v_i:i\text{ odd}\},\qquad p_s=\varnothing\ (s\notin\{a,b\}).
\]
The greatest $a$-capture profile attainable while retaining that
positivity is the disjunction of the odd variables.
\end{proposition}
\begin{proof}
Neutral elimination makes $C_m$ nonempty and core; neutral comparison gives
$T_m\subseteq C_m\cup X$. Its own-$a$ helper profiles are the odd
projections, and its source box is positive at $p$, since every helper
there is its gate's opposite atom.

Put $U=\{c,d\}$. For adjacent indices $i,i+1$, use the row
\[
 X=U,\quad v_i=\{r_i\},\quad v_{i+1}=\{r_{i+1}\},
 \quad v_j=U\ (j\notin\{i,i+1\}).
\]
At stage $i+1$ the unmarked singleton resets both the original and
neutral computations to $r_{i+1}$. Later helpers are $U$, so
$C_m\cup X$ omits $r_i$ on this row.
Only $a,b$ occur as atomic input owners. Therefore no helper can be
a singleton of another owner, and neither an extra-owned gate nor a
union attachment can erase an installed $a$- or $b$-winner.

The case $m=1$ is immediate; assume $m\geq3$.
The requested capture supplies an $a$-helper positive on $\{v_1\}$;
let $L_1$ be its last such location. On the row for $(1,2)$ this
helper resets the running value to $a$. Safety forces a later
$b$-helper positive on $\{v_2\}$, since that is the only way to
remove $a$. Thus its last location $L_2$ exists and $L_1<L_2$.
The row for $(2,3)$ then forces a later $a$-helper positive on
$\{v_3\}$. Induction gives
\[
 L_1<L_2<\cdots<L_m.
\]
All later singleton captures are forced by the first one. The odd
locations are $a$-owned. Their helper profiles are distinct: a later
one positive on an earlier odd singleton would contradict the
definition of its last location. This proves both lower bounds,
and $T_m$ attains the depth bound.

Finally, positivity at $p$ forbids capture on every subset of the
even variables. The disjunction of the odd variables is the greatest
monotone profile with that restriction, and $T_m$ attains it.
\end{proof}

By \cref{core:prop:principal-activation-boxes}, the least interval
predicate positive at this $p$ is one actual box. Its $a$-capture
profile is exactly the displayed maximum. Every one-$X$ spine attaining
this principal box therefore has at least $m$ marked-path choice gates.
Principal boxes can thus require arbitrarily many ordered helpers,
already at fixed $K=4$; even one prescribed capture bit forces this depth.
